\chapter{Équations de la droite et du plan}
\label{workbook:chapter:10}
\section{Droite dans le plan}
\label{workbook:section:10.01}
\begin{exerciseblock}{Droite dans le plan}
\item Écrire les équations vectorielle et paramétriques de la droite passant par $A(2,-1)$ et de direction $(3,4)$.
\item Trouver l'équation de la droite passant par $A(-1,2)$ et $B(5,-4)$ sous forme paramétrique puis cartésienne.
\item Déterminer si $P(8,7)$ appartient à la droite $(x,y)=(2,-1)+t(3,4)$.
\end{exerciseblock}
\section{Vecteur normal}
\label{workbook:section:10.02}
\begin{exerciseblock}{Vecteur normal}
\item Donner un vecteur normal et un vecteur directeur de la droite $3x-2y+7=0$.
\item Écrire l'équation cartésienne de la droite passant par $(4,1)$ et normale à $(2,-5)$.
\item Montrer que la direction $(b,-a)$ est orthogonale au normal $(a,b)$.
\end{exerciseblock}
\section{Positions relatives de deux droites}
\label{workbook:section:10.03}
\begin{exerciseblock}{Positions relatives de deux droites}
\item Comparer $d_1:(x,y)=(1,2)+t(2,-1)$ et $d_2:(x,y)=(5,0)+s(-4,2)$.
\item Trouver l'intersection de $2x+y=7$ et $x-y=2$.
\item Donner un critère en termes de vecteurs directeurs pour des droites parallèles, confondues ou sécantes.
\end{exerciseblock}
\section{Plan dans \texorpdfstring{$\R^3$}{R3}}
\label{workbook:section:10.04}
\begin{exerciseblock}{Plan dans \texorpdfstring{$\R^3$}{R3}}
\item Écrire l'équation du plan passant par $A(1,2,-1)$ et de vecteur normal $(2,-3,4)$.
\item Vérifier si $P(3,0,1)$ appartient à ce plan.
\item Donner une représentation paramétrique du plan passant par l'origine et engendré par $u=(1,0,2)$ et $v=(0,1,-1)$.
\end{exerciseblock}
\section{Intersection d'une droite et d'un plan}
\label{workbook:section:10.05}
\begin{exerciseblock}{Intersection d'une droite et d'un plan}
\item Trouver l'intersection de $(x,y,z)=(1,0,2)+t(2,1,-1)$ avec le plan $x+y+z=6$.
\item Déterminer les cas possibles selon que le vecteur directeur de la droite est ou non orthogonal au normal du plan.
\item Construire un exemple d'une droite contenue dans un plan et vérifier l'appartenance de deux points.
\end{exerciseblock}
\section{Point, direction et vecteur normal}
\label{workbook:section:10.06}
\begin{exerciseblock}{Point, direction et vecteur normal}
\item Pour une droite, expliquer quelles données minimales déterminent une représentation paramétrique et une représentation cartésienne.
\item Passer de $x=1+2t$, $y=3-5t$ à une équation cartésienne.
\item Passer du plan $2x-y+3z=7$ à une description par un point et deux directions indépendantes.
\end{exerciseblock}
\section{Choisir et exploiter une représentation}
\label{workbook:section:10.07}
\begin{exerciseblock}{Choisir et exploiter une représentation}
\item Pour tester l'appartenance d'un point, comparer l'efficacité des formes paramétrique et cartésienne.
\item Pour trouver une intersection, choisir une représentation adaptée entre deux droites données sous formes différentes.
\item Écrire une courte stratégie : données disponibles $\to$ représentation choisie $\to$ opération effectuée.
\end{exerciseblock}
\section*{Synthèse du chapitre}
\begin{synthesisbox}
\item Déterminer l'équation du plan passant par $A(1,0,2)$, $B(3,1,0)$ et $C(0,2,1)$.
\item Trouver la droite d'intersection des plans $x+y+z=3$ et $2x-y+z=1$.
\item Étudier la position relative de la droite $(x,y,z)=(0,1,2)+t(1,-1,2)$ et du plan $2x+y-z=4$.
\end{synthesisbox}
